\section{Version control, agents and plugins}
\label{sec:integrations}

\subsection{Git}

The \ui{Git} panel covers the day-to-day loop without leaving the editor: it
shows the current branch, how far ahead or behind you are, and your staged,
modified and untracked files. You can stage everything, write a message and
commit, push, pull, and view the diff of the file you are editing. If the
folder is not a repository yet, the panel offers to initialise one.

It is deliberately not a full Git client. Rebases, conflict resolution and
history surgery belong in a terminal or a dedicated tool. What is here is what
you want between paragraphs.

A \LaTeX{}-specific note: commit your sources, not your output. A
\texttt{.gitignore} containing

\begin{lstlisting}
build/
*.aux
*.log
*.out
*.toc
*.synctex.gz
\end{lstlisting}

keeps regenerable files out of your history. Since the editor writes all output
to \texttt{build/}, the first line does most of the work.

\subsection{The MCP server}

The editor runs a small Model Context Protocol server on
\texttt{127.0.0.1:9876}. Its purpose is to let an external coding agent see and
act on the project you have open --- useful when you want something more
open-ended than a single skill, such as restructuring a document across several
files.

It exposes six tools:

\begin{table}[htbp]
  \centering
  \caption{MCP tools}
  \label{tab:mcp}
  \begin{tabular}{ll}
    \toprule
    \textbf{Tool} & \textbf{Purpose} \\
    \midrule
    \texttt{read\_file}             & read a file from the project \\
    \texttt{write\_file}            & write or overwrite a file \\
    \texttt{compile}                & typeset and return errors and warnings \\
    \texttt{get\_errors}            & parse a \texttt{.log} for its errors \\
    \texttt{get\_project\_structure} & return the file tree \\
    \texttt{get\_section\_structure} & extract headings from a file \\
    \bottomrule
  \end{tabular}
\end{table}

The indicator at the top of the \ui{AI} panel reports whether the server
actually bound its port. If it reads \emph{not running}, something else is
using 9876; nothing else in the editor is affected.

\subsubsection*{Connecting an agent}

The server speaks JSON-RPC over HTTP. Any MCP client can point at it; for
Claude Code, add it as an HTTP server:

\begin{lstlisting}
claude mcp add --transport http latex http://127.0.0.1:9876/mcp
\end{lstlisting}

To check it by hand before wiring anything up, ask it what it can do:

\begin{lstlisting}
curl -s http://127.0.0.1:9876/mcp \
  -H 'Content-Type: application/json' \
  -d '{"jsonrpc":"2.0","id":1,"method":"tools/list"}'
\end{lstlisting}

The project the editor currently has open is what the tools operate on --- the
app pushes the open project and active file to the server as you work, so an
agent always sees what you see.

\subsubsection*{Skills or agent?}

The two AI surfaces are for different shapes of work, and it is worth being
deliberate about which you reach for.

\begin{table}[htbp]
  \centering
  \caption{Choosing between the panel and an external agent}
  \label{tab:ai-surfaces}
  \begin{tabular}{p{0.44\textwidth}p{0.44\textwidth}}
    \toprule
    \textbf{AI Skills panel} & \textbf{External agent over MCP} \\
    \midrule
    One passage, one task, one diff you approve.
      & Many files, several steps, its own plan. \\[3pt]
    You see exactly what is sent and exactly what changes.
      & It reads and writes on its own; review with \texttt{git diff}. \\[3pt]
    Right for: proofreading, fixing an error, generating a table.
      & Right for: restructuring a document, renaming a label everywhere,
        splitting a chapter. \\
    \bottomrule
  \end{tabular}
\end{table}

Because \texttt{write\_file} overwrites without asking, commit before letting an
agent loose --- then the review is a diff rather than an act of faith.

Two cautions. The server is bound to localhost and unauthenticated: anything
able to reach that port can read and write the open project, so do not forward
it off the machine. And \texttt{write\_file} overwrites without asking, so
prefer to have an agent work on a project that is committed to Git.

\subsection{Plugins}

Plugins are loaded from \texttt{\textasciitilde/.latex-editor/plugins/}, each a
folder with a \texttt{plugin.json} manifest and an entry point. The
\ui{Plugins} dialog lists what is installed and lets you inspect each one's
manifest and source.

Being straight about the current state: plugins are \emph{listed and
inspectable, but not yet executed}. The manager is a viewer. Two samples ship
in the repository's \texttt{plugins/} folder to show the intended manifest
shape --- treat the system as a preview of an interface rather than a working
extension point.

\subsection{SyncTeX}

The backend can map a source line to a PDF page and back. It is not yet wired
to a click in the preview, so there is no jump-to-source today. The groundwork
is done; the gesture is not.
